\section{Discussion and Limitations}
\label{sec:limitations}

\paragraph{Substrate displacement.}
The central pattern in the evidence is substrate displacement.  Verification
becomes stronger when evidence is moved to a stronger substrate, but stronger
substrates often travel less naturally with ordinary LLM outputs.  Metadata can
be strong if preserved, but it is fragile under copy-paste or platform
stripping.  Protocol transcripts and proof objects can be deterministic, but
they require participation and a committed statement.  Provider-side traces can
recover codewords under authorized event access, but they do not travel with
public final text.  Natural final text is portable and non-cooperative, but the
current public predicates either fail codeword recovery or become
source-mismatch spoofing targets.

\paragraph{No universal impossibility theorem.}
The paper does not prove that every conceivable public final-text method must
fail.  It reports a conditional empirical and diagnostic result: in the tested
artifacts, provider-side first-divergence traces carry recoverable evidence,
while public final-text predicates recover zero codeword blocks and expose
spoofable accept regions.  Stronger public final-text schemes would need to
show deterministic recovery and spoof resistance under their own declared
substrate and threat model.

\paragraph{P2 and P3 are surrogates, not exhaustive classes.}
The shallow predicates used in the public final-text experiments are surrogate
attempts to recover or approximate protected signal from final text.  They are
useful because they occupy the diagnostic failure mode of interest: no
codeword recovery, but a nonzero accept region.  They are not presented as an
exhaustive class of all possible public final-text verifiers.

\paragraph{Trace-bound success is not deployment proof.}
The trace-bound diagnostic measures event-access controllability.  It should
not be interpreted as public text-only verification success, natural evidence
success, phrase-decoder success, signed provenance, cryptographic provenance,
or an ownership proof.  A deployable cooperative trace-bundle system would also
need trace binding, key management, audit policy, and release engineering that
are outside the present evidence scope.

\paragraph{Attack naturalness is not yet established.}
The guided rewrite/graft attacks show public-predicate optimizability.  The
current artifact-level result establishes neither semantic preservation,
human-naturalness, nor indistinguishability from ordinary answers.  A proxy
readability audit over the guided rewrite/graft examples found 0/60 passing
rows, mostly due to incomplete sentence endings, isolated single-letter
fragments, and known broken-graft markers.  Adding a naturalness or
semantic-preservation filter could strengthen the attack, but it is not
required to show that a predicate with zero codeword recovery and
source-mismatch accepts is not a deterministic verifier.

\paragraph{Historical failures remain failures.}
Earlier routes failed for important reasons: passive evidence lacked enough
observable symbols, step-labeled outputs exposed brittle structure, phrase
decoders failed under scrubbed final text, and natural-output diagnostics
surfaced duplicate or forbidden-quality failures.  The VSG framing does not
reclassify those artifacts as positives.  It explains why the paper should
separate trace-bound controllability from public final-text verification.

\begin{table}[t]
\centering
\scriptsize
\setlength{\tabcolsep}{3pt}
\renewcommand{\arraystretch}{1.12}
\begin{tabular}{p{0.22\linewidth}p{0.19\linewidth}p{0.26\linewidth}p{0.22\linewidth}}
\toprule
Route & Status & Bottleneck & Use in VSG \\
\midrule
v1 passive mining & Historical failure & Public-symbol survival and observability failed & Motivates substrate-gap framing only \\
R3 Step-label/fixed-slot route & Historical failure & Visible structure and slot fragility & Not reclassified as positive \\
R4 phrase/final-text route & Historical failure & Public final-text codeword recovery remains zero & Supports observability gap \\
R4 trace-bound route & Diagnostic success under trace access & Not public final-text verification & Upper-bound controllability evidence only \\
\bottomrule
\end{tabular}
\caption{Historical failure chain motivating the substrate-gap framing.  The
chain does not imply public final-text verification success or ownership
proof.}
\label{tab:historical-chain}
\end{table}

\paragraph{Unsupported axes.}
Unsupported axes include sanitizer robustness, robustness to semantic
paraphrase, full false-accept-rate guarantees, payload-diversity evidence,
Llama public-text transfer, model-family general verification, copied-text
ownership resolution, rewritten-text ownership resolution, and non-cooperative
ownership proof.  These are future work or out-of-scope axes, not hidden
successes.
